\begin{abstract}
The proximal point algorithm of Martinet and Rockafellar was developed for monotone operators. Later local analyses, including that of Luke and Tam, combine weaker monotonicity assumptions with residual regularity to obtain convergence and describe the geometry of operator values and zero sets. Recent estimates of Alcantara further clarify this quantitative relation. We study how much these two properties can compensate for one another, and how global operator conditions can establish the existence of proximal points. For the ordinary set-valued proximal point method in finite dimensions, DD1 gives a piecewise compensation formula in both directions. Under local existence and localization assumptions, it yields geometric distance decay, finite length, and point convergence, with a parameter region larger than the displayed sufficient conditions used for comparison. DD2 treats general residual functions and signed inner product bounds. It guarantees a prescribed distance decrease and controls the actual step. Summability or integrability conditions then give finite length, including for sublinear trajectories; an example exhibits point convergence with infinite length. DD3 uses the homotopy and topological tools of Brito, Lara, and Yuan to prove global proximal solvability from an anchored condition uniform over a parameter interval, and combines this existence result with the local convergence theorems. Finally, we determine admissible proximal parameters from data on the original operator and prove convergence with finite length when the parameter varies arbitrarily within a common admissible range.
\end{abstract}
